\documentclass[12pt]{article}
\usepackage[utf8]{inputenc}
\usepackage[T1]{fontenc}
\usepackage{hyperref}
\title{Saya Language Resources for K-12 Curriculum Development}
\author{Compiled by Claude Sonnet 5.5 (claude-sonnet-5-5) for Andr\'as Kornai}
\date{}
\begin{document}
\maketitle

\section{Introduction}
Saya (ISO 639-3 \texttt{say}; Glottocode saya1246) is a West Chadic (Afro-Asiatic) language of southern Bauchi State, Nigeria, mainly in the Tafawa Balewa and Bogoro local government areas. It is also called Sayanci, Sayawa, Seya or Seyawa. Zaar and Sahng are treated as dialects. Wikipedia gives about 300,000 speakers (2013), while Endangered Languages cites about 150,000 and rates the language as vulnerable. The sources I found describe speakers as bilingual in Hausa and say children are schooled in Hausa and then English. I found no evidence of a Saya-medium school curriculum, and very little published educational material.

\section{Primary Educational Resources}
I found no primers, textbooks or school programmes in Saya in search results. Joshua Project lists Bible portions in Zaar (Sayawa) but no New Testament and no complete Bible. \url{https://www.joshuaproject.net/index.php/languages/say}

\section{Linguistic and Reference Materials}
\begin{itemize}
\item Wikipedia overview of the Saya language, with names, dialects, location and speaker figures: \url{https://en.wikipedia.org/wiki/Saya_language}
\item Endangered Languages entry. It reports no official orthography and cites Caron's work: \url{https://endangeredlanguages.com/lang/10434}
\item A recent proposal for Zaar orthography and spelling reform (Journal of Language and Literature Studies, JOLAN), proposing 55 symbols or patterns for 56 sounds and focused mainly on the Moss dialect: \url{https://jolan.com.ng/index.php/home/article/view/420}
\item Glottolog language page: \url{https://glottolog.org/resource/languoid/id/saya1246.bigmap.html}
\item Numeral data for Saya, from the Max Planck Institute for Evolutionary Anthropology (channumerals) database: \url{https://lingweb.eva.mpg.de/channumerals/Saya.htm}
\item South Bauchi languages (classification context): \url{https://en.wikipedia.org/wiki/South_Bauchi_languages}
\end{itemize}

\section{Religious and Community Resources}
The Zaar are largely Christian and, according to Caron as quoted by Endangered Languages, use a Hausa Bible. No Zaar-language complete scripture was found. See the Joshua Project and Endangered Languages pages above.

\section{Digital and Audio Resources}
I found no audio, video, corpus or software resources specific to Saya in my searches.

\section{Curriculum Development Context}
A curriculum would start from an orthography that is still unsettled. Caron (2005) attempted one, and the JOLAN proposal is a recent revision. Hausa is the existing school language of literacy, followed by English. Any Saya curriculum would need new primers, an agreed spelling and a decision about which dialect to standardise. Native-speaker and community consultation would be necessary because published materials are scarce.
\end{document}